\chapter{Dermatitis Herpetiformis}
\index{Dermatitis Herpetiformis}

\section*{Definition and Overview}
Dermatitis herpetiformis is a chronic, intensely itchy skin condition caused by an immune reaction to gluten, the protein found in wheat, barley, and rye. It is the skin manifestation of coeliac disease: people with dermatitis herpetiformis have the same gluten sensitivity that damages the gut, even when they have no digestive symptoms. The condition causes groups of small, red, intensely itchy blisters, typically on the elbows, knees, buttocks, and scalp. It is not related to the herpes virus; the name refers to the appearance of the blisters. Strict avoidance of gluten is the main treatment.

\section*{Epidemiology}
Dermatitis herpetiformis is uncommon, affecting around 10--30 people per 100,000, and it is more common in people of northern European descent. It typically appears between the ages of 20 and 60 but can occur at any age, and it is slightly more common in men. Almost all people with dermatitis herpetiformis have coeliac-type changes in the small intestine, and a family history of coeliac disease or dermatitis herpetiformis increases risk. Many people with the condition have no gut symptoms at all.

\section*{Aetiology and Pathophysiology}
\begin{itemize}
    \item \textbf{Gluten sensitivity:} the immune system reacts to gluten, a protein in wheat, barley, and rye
    \item \textbf{Antibody deposition:} gluten triggers the production of IgA antibodies that deposit in the skin, particularly in the dermal papillae
    \item \textbf{Neutrophil attack:} the deposited antibodies attract white blood cells that damage the skin, causing blisters
    \item \textbf{Genetic susceptibility:} the condition is strongly associated with specific HLA genes (HLA-DQ2 and HLA-DQ8), as in coeliac disease
    \item \textbf{Gut involvement:} most people have coeliac-type damage to the intestine, whether or not they have digestive symptoms
    \item \textbf{Associated conditions:} thyroid disease and other autoimmune conditions are more common
\end{itemize}

\section*{Clinical Features}
\begin{itemize}
    \item Groups of small, red, intensely itchy blisters and bumps
    \item Symmetrical distribution on the elbows, knees, buttocks, scalp, and back
    \item Intense burning and itching, often worse at night
    \item Blisters that break easily, leaving small sores and crusts
    \item Persistent scratching causes thickening and scarring of the skin
    \item Most people have no digestive symptoms, although some have bloating, diarrhoea, or malabsorption
    \item Symptoms can be triggered or worsened by a high-gluten diet
\end{itemize}

\section*{Differential Diagnosis}
\begin{itemize}
    \item Atopic dermatitis (eczema), which is more widespread and does not blister
    \item Contact dermatitis and allergic reactions
    \item Scabies, an intensely itchy infestation
    \item Pemphigus and bullous pemphigoid, rare blistering autoimmune diseases
    \item Herpes simplex or other viral rashes
    \item Linear IgA disease, a related blistering condition
\end{itemize}

\section*{When to Seek Medical Care}
Anyone with an intensely itchy, blistering rash on the elbows, knees, or buttocks should see a doctor, as dermatitis herpetiformis requires specialist diagnosis and dietary treatment. Blood tests and skin biopsy confirm the diagnosis. People diagnosed with the condition should also have assessment for coeliac disease and associated conditions.

\textbf{Contact your local emergency services for:}
\begin{itemize}
    \item A widespread blistering rash with fever or feeling very unwell
    \item Signs of severe skin infection
    \item Difficulty breathing or facial swelling suggesting an allergic emergency
\end{itemize}

\section*{Diagnosis and Investigations}
\begin{itemize}
    \item \textbf{Skin biopsy:} a biopsy taken from normal skin beside a blister shows the characteristic IgA deposits (direct immunofluorescence)
    \item \textbf{Blood tests:} coeliac antibodies (tissue transglutaminase IgA) are positive in most cases
    \item \textbf{Genetic testing:} HLA testing may be used in some cases
    \item \textbf{Gut assessment:} endoscopy and small bowel biopsy confirm coeliac disease where indicated
    \item \textbf{Screening:} thyroid function and other autoimmune tests, as associated conditions are common
\end{itemize}

\section*{Management}
\begin{itemize}
    \item \textbf{Gluten-free diet:} the definitive treatment; strict, lifelong avoidance of wheat, barley, and rye
    \item \textbf{Dapsone:} medication that controls symptoms rapidly, used while the diet takes effect
    \item \textbf{Dietitian support:} specialist guidance to achieve a complete, nutritious gluten-free diet
    \item \textbf{Symptom relief:} soothing baths, emollients, and antipruritic measures
    \item \textbf{Monitoring:} regular review of skin and blood tests to confirm control
    \item \textbf{Management of associated conditions:} treatment of thyroid disease and other autoimmune disorders
    \item \textbf{Long-term follow-up:} skin improvement lags behind diet, and healing may take months
\end{itemize}

\section*{Complications}
\begin{itemize}
    \item Skin damage, scarring, and pigment changes from chronic scratching
    \item Secondary bacterial infection of the skin
    \item Coeliac disease complications, including malabsorption, anaemia, and osteoporosis
    \item Increased risk of intestinal lymphoma in untreated coeliac disease
    \item Side effects of dapsone, including anaemia, which requires monitoring
    \item Associated autoimmune conditions, including thyroid disease and diabetes
\end{itemize}

\section*{Prognosis and Outcomes}
With a strict gluten-free diet, the skin lesions of dermatitis herpetiformis gradually clear, usually within 6--24 months, and the itching resolves. The diet also protects the gut and reduces the long-term risks of coeliac disease. Dapsone provides rapid relief while the diet takes effect, and most people achieve complete control. Once controlled, the condition is manageable, and most people lead normal lives with careful adherence to the gluten-free diet.

\section*{Prevention and Self-Care}
\begin{itemize}
    \item Follow the gluten-free diet strictly, even when the skin appears clear
    \item Seek dietitian advice to identify all sources of gluten
    \item Take medications exactly as prescribed and attend monitoring
    \item Look after the skin with emollients and gentle care
    \item Attend regular medical review, including blood tests
    \item Discuss family screening, since close relatives are at increased risk
    \item Report any new symptoms, including gut symptoms or unexplained weight loss
\end{itemize}

\section*{Sources and Further Reading}
The following reputable sources provide further information for healthcare students and patients seeking reliable, current advice about dermatitis herpetiformis.

\begin{itemize}
    \item Coeliac Australia. \textit{Dermatitis herpetiformis: information and support.} https://www.coeliac.org.au/
    \item Australasian College of Dermatologists. \textit{Dermatitis herpetiformis.} https://www.dermcoll.edu.au/
    \item DermNet. \textit{Dermatitis herpetiformis: clinical information.} https://dermnetnz.org/
    \item NHS. \textit{Dermatitis herpetiformis: symptoms and treatment.} https://www.nhs.uk/conditions/dermatitis-herpetiformis/
    \item Mayo Clinic. \textit{Dermatitis herpetiformis: overview.} https://www.mayoclinic.org/
    \item National Institute of Diabetes and Digestive and Kidney Diseases. \textit{Coeliac disease and skin manifestations.} https://www.niddk.nih.gov/
\end{itemize}
